% !PRINTED 178
\weberpage{178}
%<<CONT
solute ones there also exist relative invariants, and these can be
used for the simplification of the form problem.

If $r$ is such a relative invariant, then a certain power
of $r$, whose degree $e$ is a divisor of the degree $\mu$ of $S$,
will be an absolute invariant, and $r$ is thus determined by a \gesperrt{pure
equation} in $\Omega$.

All the substitutions of $S$ by which $r$ remains unchanged
form by themselves a group $T_r$, and the group $T_r$ is, as we
have seen in \S.\,40, a normal divisor of $S$.

Further, however, we see that every function of the variables $x$
that remains unchanged under the substitutions of the group $T_r$
can be expressed rationally in terms of $r$, that is, is contained in the
field $\Omega_r$ derived from $\Omega$ by adjunction of $r$.

For according to \S.\,40 there is a substitution $E$ in $S$, and
a primitive $e$th root of unity $\varepsilon$, such that $r$ passes under $E$ into $\varepsilon r$,
and all the values that $r$ can assume:
\[
r, \varepsilon r, \varepsilon^2 r, \ldots \varepsilon^{e-1} r,
\]
are obtained by repetition of the substitution $E$. If, further,
$\tau$ denotes a function of the $x$ that admits the substitutions of the group
$T_r$, and passes under $E$ and its powers into
\[
\tau, \tau_1, \tau_2, \ldots \tau_{e-1}
\]
then all these functions likewise admit the
substitutions of $T_r$, because $T_r$ is a normal divisor of $S$. The
function
\weq{1}{(\varepsilon^{-h}, \tau) = \tau + \varepsilon^{-h}\tau_1 + \cdots + \varepsilon^{-h(e-1)}\tau_{e-1}, \qquad (h = 0, 1, 2, \ldots e-1),}
which is formed on the analogy of the \gesperrt{Lagrange} resolvents (Bd.\,I,
\S.\,164), takes on the factor $\varepsilon^h$ upon application of the substitution $E$,
and if we therefore set
\weq{2}{(\varepsilon^{-h}, \tau) = r^h \Psi_h
}
then $\Psi_h$ is an absolute invariant, and hence contained in the
field $\Omega$.

From (1) and (2), however, it follows that
\weq{3}{e\tau = \Psi + r\Psi_1 + r^2\Psi_2 + \cdots + r^{e-1}\Psi_{e-1},}
whereby the assertion is proved.

If we set
\weq{4}{\mu = e\nu,}
%CONT>>
